\documentclass[11pt,a4paper]{article}
\newcommand{\DocVersion}{1.1.0}
\newcommand{\BuildDate}{20261006}
% --------------------------------------------------
% Packages
% --------------------------------------------------
\usepackage[utf8]{inputenc}
\usepackage[T1]{fontenc}
\usepackage[english]{babel}

\usepackage{geometry}
\usepackage{xcolor}
\usepackage{graphicx}
\usepackage{enumitem}
\usepackage{hyperref}
\usepackage{fancyhdr}
\usepackage{booktabs}
\usepackage{tabularx}
\usepackage{array}
\usepackage{longtable}

% --------------------------------------------------
% Page settings
% --------------------------------------------------
\geometry{
    top=2.5cm,
    bottom=2.5cm,
    left=2.5cm,
    right=2.5cm
}

\definecolor{companyblue}{HTML}{1F4E79}
\definecolor{lightblue}{HTML}{EAF2F8}
\definecolor{warningred}{HTML}{9C0006}
\definecolor{warningbg}{HTML}{FFC7CE}

\hypersetup{
    colorlinks=true,
    linkcolor=companyblue,
    urlcolor=companyblue
}

\setlength{\parindent}{0pt}
\setlength{\parskip}{0.6em}

% --------------------------------------------------
% Header and footer
% --------------------------------------------------
\pagestyle{fancy}
\fancyhf{}

\lhead{Tex C. Lab}
\rhead{User Manual}
\cfoot{\thepage}

% --------------------------------------------------
% Custom commands
% --------------------------------------------------
\newcommand{\software}{\textbf{Tex C. Lab}}
\newcommand{\button}[1]{\textbf{#1}}
\newcommand{\filename}[1]{\texttt{#1}}

% --------------------------------------------------
% Document
% --------------------------------------------------
\begin{document}

% ==================================================
% Cover page
% ==================================================
\begin{titlepage}
    \centering

    \vspace*{1.5cm}

    % Company logo
    % The file must be located in Docs/imgs/logo.png
    \includegraphics[width=0.35\textwidth]{imgs/cra-logo.png}

    \vspace{1.5cm}

    {\Huge\bfseries Tex C. Lab\par}

    \vspace{0.5cm}

    {\Large Utility for managing encrypted files for Power systems\par}

    \vspace{2cm}

    {\LARGE\bfseries User Manual\par}

    \vspace{1cm}

    {\large Version \DocVersion\par}

    \vspace{0.3cm}

    {\large Build date: \BuildDate\par}

    \vfill

    {\large Tex \par}

\end{titlepage}

% ==================================================
% Table of contents
% ==================================================
\tableofcontents
\newpage

% ==================================================
\section{Introduction}
% ==================================================

\software{} is a software tool developed to generate encrypted
files used by Power systems and their Boot Manager.

The program can generate the files required to:

\begin{itemize}
    \item Automatically log in to the Boot Manager;
    \item Change the Administrator user's password;
    \item Reset the Administrator user's password;
    \item Update the card parameters;
    \item Validate new firmware for the card.
\end{itemize}

The generated files must be transferred to the card using
a USB device.

This manual describes the features available in the application
and the procedures to follow when using the generated files with
the Boot Manager.

% ==================================================
\section{Requirements and general operation}
% ==================================================

The Boot Manager manages two access levels:

\begin{itemize}
    \item \textbf{Guest};
    \item \textbf{Administrator}.
\end{itemize}

Automatic login and operations related to the Administrator user's
password are managed through files generated by \software.

To process these files, it is necessary to:

\begin{enumerate}
    \item Copy the generated file to a USB device;
    \item Connect the USB device to the card;
    \item Power on or restart the card;
    \item Press and hold the \texttt{X} button during startup;
    \item Enter the \textbf{pre-bootloader};
    \item Select the required operation.
\end{enumerate}

\subsection{Pre-bootloader}

The pre-bootloader is the startup environment accessed by holding
the \texttt{X} button while powering on the card.

From this environment, it is possible to:

\begin{itemize}
    \item Access the Boot Manager main menu;
    \item Load the card parameters;
    \item Load new firmware.
\end{itemize}

To process the generated files, it is necessary to go through the
pre-bootloader. For login, change, and reset operations, it is
sufficient to return to the main menu. For parameters and firmware,
the corresponding options must instead be selected.

% ==================================================
\section{Application interface}
% ==================================================

The application consists of the following pages:

\begin{itemize}
    \item \textbf{Manual};
    \item \textbf{About};
    \item \textbf{Settings};
    \item \textbf{Generators}.
\end{itemize}

% ==================================================
\subsection{Manual page}
% ==================================================

The Manual page redirects to a PDF viewer and automatically opens
the user manual.

% ==================================================
\subsection{About page}
% ==================================================

The \textbf{About} page contains information about the application,
including:

\begin{itemize}
    \item Software name;
    \item Installed version;
    \item Information about the software manufacturer.
\end{itemize}

% ==================================================
\subsection{Settings page}
% ==================================================

The \textbf{Settings} page allows you to select the destination
directory for files generated by the application.

\subsubsection*{Selecting the output directory}

To set the destination directory:

\begin{enumerate}
    \item Open the \button{Settings} page;
    \item Select the desired directory;
    \item Confirm the selection.
\end{enumerate}

All subsequently generated files will be saved in the configured
directory.

% ==================================================
\subsection{Generators page}
% ==================================================

The \textbf{Generators} page contains the generators used to prepare
files for use with the Boot Manager.

The available generators are:

\begin{itemize}
    \item \button{Generate Login};
    \item \button{Change Password};
    \item \button{Reset Password};
    \item \button{Card Parameters};
    \item \button{Firmware Checksum}.
\end{itemize}

% ==================================================
\section{Managing generated files}
% ==================================================

The files generated by \software{} are divided into two categories.

\subsection{Automatically processed files}

The following files are automatically processed by the pre-bootloader
when they are present on the USB device:

\begin{itemize}
    \item Login file;
    \item Password change file;
    \item Password reset file.
\end{itemize}

For this reason, these files must be removed from the USB device
after use, when they are no longer needed.

\textbf{Warning:} keeping these files on a USB device connected to
the card could cause the associated operation to be executed
automatically during a subsequent pre-bootloader startup.

\subsection{Manually loaded files}

The following files are manually selected by the user directly
from the pre-bootloader:

\begin{itemize}
    \item Card parameter file;
    \item Firmware file.
\end{itemize}

These files are not processed automatically and can therefore remain
on the USB device without causing the operation to be executed
unintentionally.

% ==================================================
\section{Generators}

The generators are designed to produce files that are automatically
processed by the bootloader, \textbf{which is accessed by pressing
and holding the controller button}.

% ==================================================

\subsection{Generate Login}

The \button{Generate Login} generator creates a file that allows
automatic login to the Boot Manager using the Administrator
credentials.

\subsubsection*{Required data}

\begin{itemize}
    \item Current Administrator password.
\end{itemize}

\subsubsection*{Procedure}

\begin{enumerate}
    \item Open the \button{Generators} page.
    \item Select \button{Generate Login}.
    \item Enter the Administrator password.
    \item Start generating the file.
    \item Copy the generated file to a USB device.
    \item Connect the USB device to the card.
    \item Power on or restart the card while holding down
    the \texttt{X} button.
    \item Enter the pre-bootloader.
    \item Select access to the \textbf{Main Menu}.
    \item Allow the Boot Manager to process the login file.
\end{enumerate}

% ==================================================
\subsection{Change Password}

The \button{Change Password} generator creates a file that allows
the Administrator password to be changed.

To perform this operation, both the current login file and the new
password are required.

\subsubsection*{Required data}

\begin{itemize}
    \item New Administrator password.
\end{itemize}

\subsubsection*{Procedure}

\begin{enumerate}
    \item Open the \button{Generators} page.
    \item Select \button{Change Password}.
    \item Select the current login file.
    \item Enter the new password.
    \item Start generating the file.
    \item Copy the generated file to a USB device.
    \item ALSO copy the current login file.
    \item Connect the USB device to the card.
    \item Power on or restart the card while holding down
    the \texttt{X} button.
    \item Enter the pre-bootloader.
    \item Select the \textbf{Main Menu}.
    \item Allow the Boot Manager to process the generated file.
\end{enumerate}

After use, remove both the login file and the password change file
from the USB device. If you want to continue using a login file,
you must regenerate it using the new password that was just set.

% ==================================================
\subsection{Reset Password}

The \button{Reset Password} generator creates a temporary file
for resetting the Administrator password.

The reset file is sufficient on its own: it is not necessary to add
a login file or a password change file.

\subsubsection*{Required data}

\begin{itemize}
    \item MAC address of the card;
    \item Expiration date and time of the file.
\end{itemize}

The MAC address must be entered in one of the formats supported
by the application, for example:

\begin{itemize}
    \item \texttt{00:11:22:33:44:55};
    \item \texttt{00-11-22-33-44-55};
    \item \texttt{001122334455}
\end{itemize}

The expiration date and time must be set in the future. The file
contains a Unix timestamp used by the Boot Manager to verify
its validity.

\subsubsection*{Procedure}

\begin{enumerate}
    \item Open the \button{Generators} page.
    \item Select \button{Reset Password}.
    \item Enter the correct MAC address of the card.
    \item Set a future expiration date and time.
    \item Start generating the file.
    \item Copy the generated file to a USB device.
    \item Connect the USB device to the card.
    \item Power on or restart the card while holding down
    the \texttt{X} button.
    \item Enter the pre-bootloader.
    \item Select the \textbf{Main Menu}.
    \item Allow the Boot Manager to process the reset file.
\end{enumerate}

The file cannot be used after the specified expiration time
and must be removed from the USB device after use.

% ==================================================
\subsection{Card Parameters}
% ==================================================

The \button{Card Parameters} generator prepares the card parameter
file for loading through the pre-bootloader.

The parameters must first be downloaded from the Boot Manager.
The resulting file, named \filename{POWER.CRD}, can be edited
as a text file.

After editing, the file must be processed by \software{}.
The program generates the encrypted file:

\begin{center}
    \filename{CARDPARS.ENC}
\end{center}

\subsubsection*{Procedure}

\begin{enumerate}
    \item Download the card parameters from the Boot Manager.
    \item Open the \filename{POWER.CRD} file with a text editor.
    \item Modify the required parameters.
    \item Save the modified file.
    \item Open the \button{Generators} page.
    \item Select \button{Card Parameters}.
    \item Select the \filename{POWER.CRD} file.
    \item Start generating the file.
    \item Copy \filename{CARDPARS.ENC} to a USB device.
    \item Connect the USB device to the card.
    \item Power on or restart the card while holding down
    the \texttt{X} button.
    \item From the pre-bootloader, select
    \button{Upload Card Parameters}.
    \item Select or confirm the \filename{CARDPARS.ENC} file.
    \item Wait for the operation to complete.
\end{enumerate}

The \filename{CARDPARS.ENC} file is manually selected by the
pre-bootloader and is not processed automatically. Therefore,
if desired, it can be left on the USB device.

% ==================================================
\subsection{Firmware Checksum}
% ==================================================

The \button{Firmware Checksum} generator validates new firmware
for loading onto the card by generating the checksum required
by the card.

The firmware must therefore first be processed by \software{}
before it can be loaded by the Boot Manager.

The program generates the encrypted file:

\begin{center}
    \filename{CHECKSUM.ENC}
\end{center}

\subsubsection*{Required data}

\begin{itemize}
    \item Firmware binary file.
\end{itemize}

\subsubsection*{Procedure}

\begin{enumerate}
    \item Open the \button{Generators} page.
    \item Select \button{Firmware Checksum}.
    \item Select the firmware binary file.
    \item Start generating the file.
    \item Copy \filename{CHECKSUM.ENC} to a USB device
    (together with the firmware itself).
    \item Connect the USB device to the card.
    \item Power on or restart the card while holding down
    the \texttt{X} button.
    \item From the pre-bootloader, select
    \button{Upload Firmware}.
    \item Wait for the update to complete.
\end{enumerate}

The \filename{CHECKSUM.ENC} file is manually selected by the
pre-bootloader and is not processed automatically.

% ==================================================
\section{Generator summary}
% ==================================================

\renewcommand{\arraystretch}{1.4}

\begin{longtable}{p{3.2cm} p{4.2cm} p{5.2cm}}
\toprule
\textbf{Generator} &
\textbf{Input} &
\textbf{Result and usage} \\
\midrule
\endhead

\button{Generate Login} &
Administrator password &
Automatic login file to be used from the Main Menu. \\

\button{Change Password} &
Current login file and new password &
File for changing the Administrator password from the Main Menu. \\

\button{Reset Password} &
MAC address and expiration &
Temporary file for resetting the password from the Main Menu. \\

\button{Card Parameters} &
Modified \filename{POWER.CRD} file &
Generates \filename{CARDPARS.ENC}, to be used with
\button{Upload Card Parameters}. \\

\button{Firmware Checksum} &
Firmware binary file &
Generates \filename{CHECKSUM.ENC}, to be used with
\button{Upload Firmware}. \\

\bottomrule
\end{longtable}

% ==================================================
\section{Warnings}
% ==================================================

\begin{itemize}
    \item Always verify that the USB device contains the
    correct file before connecting it to the card.

    \item Remove the password change and reset files after use,
    as they may be processed automatically by the pre-bootloader.

    \item After changing the password, if you want to use automatic
    login, remember to regenerate the login file using the new
    password that was just set.

    \item Carefully verify the MAC address before generating
    a reset file.

    \item Make sure that the reset file's expiration date and
    time are in the future.

    \item Do not manually modify encrypted files generated
    by the application.

    \item Do not interrupt power to the card during parameter
    or firmware updates.

    \item Use only firmware compatible with the card model.
\end{itemize}

% ==================================================
\section{Troubleshooting}
% ==================================================

\subsection{The file is not processed}

Check the following:

\begin{itemize}
    \item the file is located in the correct directory on the USB device;
    \item the file name has not been changed;
    \item the USB device was connected before startup;
    \item the \texttt{X} button was held down during power-on;
    \item the correct option was selected in the Boot Manager;
    \item the file has not expired;
    \item the file was generated for the correct card.
\end{itemize}

\subsection{The reset file is not valid}

Check that:

\begin{itemize}
    \item the entered MAC address corresponds to the card;
    \item the MAC address format is correct;
    \item the card's date and time are correct;
    \item the file has not passed its expiration time.
\end{itemize}

\subsection{Firmware upload fails}

Check that:

\begin{itemize}
    \item the selected firmware is compatible with the card;
    \item the file was processed using
    \button{Firmware Checksum};
    \item the \filename{CHECKSUM.ENC} file was used;
    \item the operation was started from the
    \button{Upload Firmware} option.
\end{itemize}

\end{document}

